% 杨昊东论文可见版式的LaTeX统一样式层
% 目标：章节、图表、公式、目录、摘要、参考文献等版式与参考论文保持同一体系。
% 说明：学校封面中的姓名、导师、专业等个体信息不在此文件中硬编码。

\usepackage{amsmath,amssymb,graphicx,booktabs,multirow}
\usepackage{geometry}
\usepackage{caption}
\usepackage{fancyhdr}
\usepackage{chngcntr}
\usepackage{setspace}
\usepackage{indentfirst}
\usepackage{titlesec}
\usepackage{tocloft}
\usepackage{hyperref}
\usepackage{etoolbox}

% A4学位论文正文版心。以参考论文的Word学位论文版式为基准统一。
\geometry{a4paper,top=2.54cm,bottom=2.54cm,left=3.00cm,right=2.60cm,headheight=14pt,headsep=0.60cm,footskip=1.20cm}

% 正文字体与段落
\zihao{-4}
\setstretch{1.5}
\setlength{\parindent}{2em}
\setlength{\parskip}{0pt}

% 将ctexart的一级section作为“章”使用：第一章、第二章……
\ctexset{
  section = {
    name = {第,章},
    number = \chinese{section},
    format = \centering\heiti\zihao{3},
    beforeskip = 0pt,
    afterskip = 1.2\baselineskip
  },
  subsection = {
    format = \heiti\zihao{-3},
    beforeskip = 0.8\baselineskip,
    afterskip = 0.4\baselineskip
  },
  subsubsection = {
    format = \heiti\zihao{4},
    beforeskip = 0.6\baselineskip,
    afterskip = 0.3\baselineskip
  }
}

% 每章另起一页；首章也如此
\pretocmd{\section}{\clearpage}{}{}

% 节标题按1.1、1.1.1编号
\renewcommand{\thesubsection}{\arabic{section}.\arabic{subsection}}
\renewcommand{\thesubsubsection}{\arabic{section}.\arabic{subsection}.\arabic{subsubsection}}

% 图、表、公式按章编号：图2-1、表2-1、式(2-1)
\counterwithin{figure}{section}
\counterwithin{table}{section}
\numberwithin{equation}{section}
\renewcommand{\thefigure}{\arabic{section}-\arabic{figure}}
\renewcommand{\thetable}{\arabic{section}-\arabic{table}}
\renewcommand{\theequation}{\arabic{section}-\arabic{equation}}
\renewcommand{\figurename}{图}
\renewcommand{\tablename}{表}

% 图表题注居中，题注字号统一
\captionsetup{font={small},labelsep=space,justification=centering,singlelinecheck=true}
\captionsetup[table]{position=top,skip=6pt}
\captionsetup[figure]{position=bottom,skip=6pt}

% 页眉页脚：正文保持学位论文简洁形式
\pagestyle{fancy}
\fancyhf{}
\fancyfoot[C]{\thepage}
\renewcommand{\headrulewidth}{0pt}
\renewcommand{\footrulewidth}{0pt}

% 目录：列至三级标题，与参考论文目录层级一致
\setcounter{tocdepth}{3}
\renewcommand{\contentsname}{目录}
\renewcommand{\cftsecfont}{\songti}
\renewcommand{\cftsubsecfont}{\songti}
\renewcommand{\cftsubsubsecfont}{\songti}
\renewcommand{\cftsecpagefont}{\songti}
\renewcommand{\cftsubsecpagefont}{\songti}
\renewcommand{\cftsubsubsecpagefont}{\songti}

% 超链接仅保留PDF内部跳转，不出现彩色框线
\hypersetup{hidelinks,unicode=true}

% 摘要页标题命令
\newcommand{\cnabstracttitle}{%
  \begin{center}\heiti\zihao{3} 摘\quad 要\end{center}\vspace{0.8\baselineskip}}
\newcommand{\enabstracttitle}{%
  \begin{center}\bfseries\zihao{3} Abstract\end{center}\vspace{0.8\baselineskip}}

% 关键词格式
\newcommand{\cnkeywords}[1]{\par\vspace{0.8\baselineskip}\noindent\textbf{关键词：}#1}
\newcommand{\enkeywords}[1]{\par\vspace{0.8\baselineskip}\noindent\textbf{Keywords: }#1}

% 参考文献标题和目录入口
\renewcommand{\refname}{参考文献}
\newcommand{\printthesisbibliography}{%
  \clearpage
  \phantomsection
  \addcontentsline{toc}{section}{参考文献}
  \bibliographystyle{gbt7714-numerical}
  \bibliography{references/publication,references/methods,references/local_theses,references/kg_sources}
}
